% Unofficial Purdue University Poster Template
% based on
% https://github.com/k4rtik/uchicago-poster
% a fork of https://github.com/anishathalye/gemini

\documentclass[final]{beamer}

% ====================
% Packages
% ====================

\usepackage[T1]{fontenc}
\usepackage{lmodern}
\usepackage[size=custom,width=120,height=72,scale=1.0]{beamerposter}
% Self-contained Purdue-inspired theme; no external Gemini files required.
\definecolor{PurdueGold}{HTML}{CFB991}
\definecolor{PurdueBlack}{HTML}{000000}
\definecolor{PurdueLight}{HTML}{F5F1E9}
\setbeamercolor{normal text}{fg=PurdueBlack,bg=white}
\setbeamercolor{headline}{fg=white,bg=PurdueBlack}
\setbeamercolor{footline}{fg=PurdueBlack,bg=PurdueGold}
\setbeamercolor{block title}{fg=PurdueBlack,bg=white}
\setbeamercolor{block body}{fg=PurdueBlack,bg=white}
\setbeamercolor{block title alerted}{fg=PurdueBlack,bg=PurdueGold}
\setbeamercolor{block body alerted}{fg=PurdueBlack,bg=PurdueLight}
\setbeamercolor{block title example}{fg=PurdueBlack,bg=PurdueGold}
\setbeamercolor{block body example}{fg=PurdueBlack,bg=PurdueLight}
\setbeamercolor{itemize item}{fg=PurdueBlack}
\setbeamercolor{itemize subitem}{fg=PurdueBlack}
\setbeamercolor{enumerate item}{fg=PurdueBlack}
\setbeamerfont{title}{size=\fontsize{64}{72}\selectfont,series=\bfseries}
\setbeamerfont{author}{size=\fontsize{32}{40}\selectfont}
\setbeamerfont{institute}{size=\fontsize{28}{36}\selectfont}
\setbeamerfont{block title}{size=\fontsize{34}{40}\selectfont,series=\bfseries}
\setbeamerfont{block title alerted}{parent=block title}
\setbeamerfont{block title example}{parent=block title}
\setbeamertemplate{navigation symbols}{}
\setbeamertemplate{caption}[numbered]
\setbeamertemplate{blocks}[default]
\newcommand{\heading}[1]{\par\medskip{\bfseries #1}\par\smallskip}
\newcommand{\footercontent}[1]{\def\posterfooter{#1}}
\setbeamertemplate{headline}{%
  \begin{beamercolorbox}[wd=\paperwidth,sep=1.5cm,center]{headline}
    {\includegraphics[width=23cm]{images/purdue-logo.png}\par}
    \vskip0.6cm
    {\usebeamerfont{title}\inserttitle\par}
    \vskip0.6cm
    {\usebeamerfont{author}\insertauthor\par}
    \vskip0.4cm
    {\usebeamerfont{institute}\insertinstitute\par}
  \end{beamercolorbox}
  {\color{PurdueGold}\rule{\paperwidth}{0.35cm}}
  \vskip0.8cm
}
\setbeamertemplate{footline}{%
  \begin{beamercolorbox}[wd=\paperwidth,sep=0.7cm]{footline}
    \fontsize{22}{28}\selectfont\posterfooter
  \end{beamercolorbox}
}
\usepackage{graphicx}
\usepackage{booktabs}
\usepackage{doi}
\usepackage[numbers]{natbib}
\usepackage[patch=none]{microtype}
\usepackage{tikz}
\usepackage{pgfplots}
\pgfplotsset{compat=1.18}
\usepackage{anyfontsize}

\pdfstringdefDisableCommands{%
\def\translate#1{#1}%
}

% ====================
% Lengths
% ====================

% If you have N columns, choose \sepwidth and \colwidth such that
% (N+1)*\sepwidth + N*\colwidth = \paperwidth
\newlength{\sepwidth}
\newlength{\colwidth}
\setlength{\sepwidth}{0.025\paperwidth}
\setlength{\colwidth}{0.3\paperwidth}

\newcommand{\separatorcolumn}{\begin{column}{\sepwidth}\end{column}}


% Draft based on user-provided project abstract and repository code at
% 87a5b301c5aca2775e9d210042a21a07fc038493 (2026-11-02).
% Five-metric demonstration is supplied by the team; no accuracy claim is implied.
% Authors supplied by the user in the Team 2 membership screenshot.
% Title follows the user-confirmed onboarding framing.
% TODO: Confirm venue, demonstration PDF and metric names.
% Cross-team roles: user-shared Landseer Agent Tools.xlsx, team 1--4 sheets.
\title{Agentic Onboarding Tools for Empirical\\Evaluation of Machine Learning Defenses}
\author{Aarush Kumar \and Andrew Zeng \and Benjamin Taylor \and Bhanu Adaveni \and Rudra Vaishnav \and Yuchao Shen}
\institute{Team 2 \enspace | \enspace Purdue University}
\footercontent{\href{https://github.com/Auspr7ng/fair_fairness_benchmark}{github.com/Auspr7ng/fair\_fairness\_benchmark}\hfill Working draft \enspace | \enspace November 2026}
\setbeamercolor{caption name}{fg=PurdueBlack}
\setbeamerfont{block body}{size=\fontsize{27}{34}\selectfont}
\setbeamerfont{block body alerted}{parent=block body}
\setbeamerfont{block body example}{parent=block body}
\setlength{\parskip}{0.45em}
\begin{document}
\begin{frame}[t]
\begin{columns}[t]
\separatorcolumn
\begin{column}{\colwidth}

\begin{block}{1 \enspace Motivation: multiple guarantees}
Real-world machine learning systems may need to provide \textbf{privacy, robustness, and fairness simultaneously}. Defenses designed for individual guarantees can interact, changing their combined effects.

To study these interactions empirically, existing defenses must be brought into a common evaluation workflow. This \textbf{onboarding} is difficult: research papers and implementations use different metrics, experimental settings, and reporting formats.

We are developing \textbf{agentic onboarding tools} for this broader effort. Our team's focus is the paper-analysis step: identifying evaluation metrics and organizing reported values into reference records for later validation.
\end{block}
\vspace{0.7cm}

\begin{alertblock}{Research question}
How can agentic onboarding connect published evaluation results to structured records for empirical ML defense evaluation?
\end{alertblock}
\vspace{0.7cm}

\begin{block}{2 \enspace Our contribution: three tools}
\begin{itemize}
\item \textbf{Metric extractor:} identify evaluation metric names from research papers.
\item \textbf{Record initializer:} structure paper-reported numerical values into metric records.
\item \textbf{Insertion tool:} validate metric names and numerical updates, with JSON storage.
\end{itemize}

We use the \textbf{Fair Fairness Benchmark (FFB)} repository as a local development and demonstration setting. The original FFB benchmark and published results are prior work~[1]; our contribution is the supporting paper-analysis workflow~[2].
\end{block}
\vspace{0.7cm}

\begin{block}{Team 2: paper-to-record workflow}
\begin{center}
\begin{tikzpicture}[x=1cm,y=1cm]
\tikzset{paperstage/.style={draw=PurdueBlack,fill=PurdueLight,line width=1.4pt,text width=29cm,align=center,inner sep=0.45cm,font=\fontsize{26}{32}\selectfont},paperflow/.style={->,line width=1.8pt}}
\node[paperstage] (pdf) at (0,0) {\textbf{Research paper PDF}};
\node[paperstage] (names) at (0,-3.1) {\textbf{Extractor}\\Identify evaluation metric names};
\node[paperstage] (values) at (0,-6.7) {\textbf{Initializer}\\Find reported values; optionally select a table row};
\node[paperstage,fill=PurdueGold] (json) at (0,-10.3) {\textbf{Insertion tool + JSON record}\\Validate and persist reviewed metric updates};
\node[paperstage] (compare) at (0,-13.9) {\textbf{Future: empirical validation}\\Compare with reproduced experimental results};
\draw[paperflow] (pdf)--(names);
\draw[paperflow] (names)--(values);
\draw[paperflow] (values)--(json);
\draw[paperflow,dashed] (json)--(compare);
\end{tikzpicture}
\end{center}
{\fontsize{22}{28}\selectfont Unresolved entries remain \texttt{null}; full dataset, model, and row provenance remains future work. The dashed final step denotes planned integration into the broader empirical framework.}
\end{block}
\end{column}
\separatorcolumn
\begin{column}{\colwidth}

\begin{block}{3 \enspace How the four teams connect}
\begin{center}
\begin{tikzpicture}[x=1cm,y=1cm]
\tikzset{team/.style={draw=PurdueBlack,fill=PurdueLight,line width=1.5pt,text width=13cm,align=center,inner sep=0.55cm,font=\fontsize{25}{31}\selectfont},handoff/.style={->,dashed,line width=1.8pt},labeltext/.style={font=\fontsize{21}{26}\selectfont,align=center,fill=white,inner sep=0.2cm}}
\node[team,text width=29cm] (input) at (0,0) {\textbf{Research paper + source repository}};
\node[team,fill=PurdueGold,line width=2.5pt] (two) at (-8,-5.6) {\textbf{TEAM 2 --- OUR FOCUS}\\Paper reference records\\[0.25cm]Extract metric names\\Initialize reported values\\Validate updates; save JSON};
\node[team] (three) at (8,-4) {\textbf{TEAM 3}\\Repository / model adaptation\\[0.2cm]Clone repositories\\Identify datasets; add models};
\node[team] (one) at (8,-9.4) {\textbf{TEAM 1}\\Environment preparation\\[0.2cm]Inspect dependencies\\Generate Dockerfiles\\Build / provision containers};
\node[team,text width=29cm] (four) at (0,-15.5) {\textbf{TEAM 4 --- REPRODUCE + VALIDATE}\\Run experiments; extract reproduced metrics\\Compare with paper reference values};
\draw[handoff] (input.south) -- ++(0,-0.8) -| (two.north);
\draw[handoff] (input.south) -- ++(0,-0.8) -| (three.north);
\draw[handoff] (three.south)--(one.north);
\draw[handoff] (two.south)--(-8,-12.7) node[labeltext,midway] {Reference\\records}--(four.north -| two.south);
\draw[handoff] (one.south)--(four.north -| one.south) node[labeltext,midway] {Execution\\environment};
\end{tikzpicture}
\end{center}
{\fontsize{22}{28}\selectfont\textbf{Dashed arrows: proposed handoffs.} Roles follow the shared tool inventory~[3]. Adaptation and environment preparation may iterate; this diagram does not claim a completed integration.}
\end{block}
\vspace{0.5cm}

\begin{block}{4 \enspace Tool design}
\heading{Extractor: paper $\rightarrow$ metric names}
\textbf{Input:} research paper PDF. \textbf{Output:} a list of evaluation metric names. Extracted text and tables feed an LLM-assisted selector, with a dictionary fallback when the model is unavailable or fails.

\heading{Initializer: names + paper $\rightarrow$ record}
\textbf{Input:} paper, optional metric names, and optional table-row label. \textbf{Output:} a metric-to-value record. Heuristics rank candidate values; unresolved entries remain \texttt{null} for review.

\heading{Insertion: reviewed value $\rightarrow$ saved record}
\textbf{Input:} existing record, metric name, and numerical value. \textbf{Output:} the updated record, with JSON loading and saving. Unknown metrics and invalid numerical values are rejected.

Local agent interfaces expose these operations, and an offline entry point connects metric selection to record initialization~[2].
\end{block}
\end{column}
\separatorcolumn
\begin{column}{\colwidth}

\begin{block}{5 \enspace Preliminary demonstration}
\begin{beamercolorbox}[sep=0.6cm,center]{block body alerted}
{\fontsize{74}{80}\selectfont\bfseries 5}\par
\textbf{evaluation metric names identified}\par
in a local FFB demonstration
\end{beamercolorbox}

This team-reported demonstration establishes preliminary \textbf{metric-extraction functionality}. It does not yet measure extraction accuracy, value correctness, or end-to-end reproduction success.

The structured records are intended to provide reference values for comparison with reproduced results. This comparison has not yet been evaluated end to end.
\end{block}
\vspace{0.6cm}

\begin{block}{6 \enspace Current status and limitations}
\begin{itemize}
\item Three tools and local agent interfaces are implemented in the repository.
\item Selected-row targeting is supported; complete dataset, model, and row provenance remains future work.
\item PDF layouts, repeated results, and heuristic value selection still require human review.
\item Agent tool-selection quality and the complete onboarding-to-validation workflow remain to be evaluated.
\end{itemize}
\end{block}
\vspace{0.6cm}

\begin{alertblock}{7 \enspace Next steps}
\begin{enumerate}
\item Build a manually annotated set of metrics, values, and source locations across papers.
\item Evaluate metric precision/recall, value correctness, and experimental-context matching.
\item Integrate the three tools into the broader empirical framework and evaluate end-to-end defense reproduction and validation.
\end{enumerate}
\end{alertblock}
\vspace{0.6cm}

\begin{block}{Takeaway}
Agentic onboarding can help connect published ML defense evaluations to structured, reviewable reference records.

Our prototype implements metric extraction, record initialization, and validated value insertion. The next milestone is to evaluate these tools within an integrated reproduction and validation workflow.
\end{block}
\vspace{0.3cm}

\begin{block}{References}
{\fontsize{21}{27}\selectfont
[1] Han et al. \emph{FFB: A Fair Fairness Benchmark for In-Processing Group Fairness Methods}. ICLR 2024. \href{https://arxiv.org/abs/2306.09468}{arXiv:2306.09468}.\par
[2] Team repository, code snapshot \texttt{87a5b30}, Nov. 2, 2026. \href{https://github.com/Auspr7ng/fair_fairness_benchmark}{Auspr7ng/fair\_fairness\_benchmark}.\par
[3] \emph{Landseer Agent Tools}, shared project workbook, team worksheets.\par
}
\end{block}
\end{column}
\separatorcolumn
\end{columns}
\end{frame}
\end{document}
